\documentclass[12pt,a4paper]{article}
\usepackage[utf8]{inputenc}
\usepackage[T1]{fontenc}
\usepackage{times}
\usepackage[margin=1in]{geometry}
\usepackage{graphicx}
\usepackage{float}
\usepackage{enumitem}
\usepackage{tabularx}
\usepackage{multirow}
\usepackage{longtable}
\usepackage{array}
\usepackage{booktabs}
\usepackage{fancyhdr}
\usepackage{lastpage}
\usepackage{hyperref}
\usepackage{tocloft}
\usepackage{listings}
\usepackage{color}
\usepackage{amsmath}
\usepackage{algorithm}
\usepackage{algorithmic}
\usepackage{tikz}
\usetikzlibrary{shapes,arrows,positioning}

% Define colors
\definecolor{lightgray}{gray}{0.95}
\definecolor{darkgray}{gray}{0.3}
\definecolor{orange}{RGB}{255,140,0}

% Page setup
\pagestyle{fancy}
\fancyhf{}
\fancyhead[L]{Software Design Document}
\fancyhead[R]{NexAgent}
\fancyfoot[C]{Page \thepage\ of \pageref{LastPage}}
\renewcommand{\headrulewidth}{0.4pt}
\renewcommand{\footrulewidth}{0.4pt}

% Section formatting
\renewcommand{\cftsecleader}{\cftdotfill{\cftdotsep}}
\renewcommand{\cfttoctitlefont}{\Large\bfseries}
\setlength{\parindent}{0pt}
\setlength{\parskip}{6pt}

% Code listing style
\lstset{
    basicstyle=\footnotesize\ttfamily,
    backgroundcolor=\color{lightgray},
    breaklines=true,
    frame=single,
    numbers=left,
    numberstyle=\tiny\color{darkgray},
    captionpos=b,
    keepspaces=true,
    showspaces=false,
    showstringspaces=false,
    showtabs=false,
    tabsize=2
}

% Document begins
\begin{document}

% Title Page
\begin{titlepage}
    \centering
    \vspace*{1cm}
    
    \rule{\linewidth}{0.5mm} \\[0.4cm]
    {\huge\bfseries Software Design Document (SDD)}\\[0.4cm]
    \rule{\linewidth}{0.5mm} \\[1.5cm]
    
    {\Large\bfseries NexAgent}\\[0.5cm]
    {\large AI Workflow Automation Platform}\\[2cm]
    
    {\Large Version 1.0}\\[2cm]
    
    {\large\bfseries By}\\[0.5cm]
    {\large
    \begin{tabular}{ll}
        Maarij Bukhari & 221803\\
        Syed M Taha & 221787\\
        Abdullah & 221837\\
    \end{tabular}
    }\\[2cm]
    
    {\large\bfseries Supervisor}\\[0.5cm]
    {\large Dr. Humaira Waqas}\\[2cm]
    
    {\large Bachelor of Science in Software Engineering}\\[0.5cm]
    {\large Air University, Islamabad}\\[1cm]
    
    {\large November 1, 2025}
    
    \vfill
\end{titlepage}

% Table of Contents
\tableofcontents
\newpage

% Revision History
\section*{Revision History}
\addcontentsline{toc}{section}{Revision History}

\begin{table}[h]
\centering
\begin{tabularx}{\textwidth}{|c|c|X|c|}
\hline
\textbf{Version} & \textbf{Date} & \textbf{Description} & \textbf{Author}\\
\hline
1.0 & November 1, 2025 & Initial SDD Document & Team\\
\hline
\end{tabularx}
\end{table}

\newpage

% Main Document Content
\section{Introduction}

\subsection{Purpose}
This Software Design Document (SDD) provides a comprehensive design overview of the NexAgent system, a no-code AI workflow automation platform. This document describes the architecture, components, interfaces, and data design of the system to guide implementation and maintenance.

\subsection{Scope}
NexAgent is a web-based platform that enables users to create, manage, and execute automated workflows without coding knowledge. The system encompasses the following 9 core modules:

\begin{enumerate}
    \item User Authentication Module
    \item Workflow Builder Module
    \item Canvas Interaction Module
    \item Workflow Execution Engine Module
    \item User Management Module
    \item Usage Analytics Module
    \item Error Handling \& Logging Module
    \item Multi-Model Integration Module
    \item Settings \& Configuration Module
\end{enumerate}

\subsection{System Verification Status}
\textbf{Research Completed:} November 1, 2025\\
\textbf{Verification Method:} Direct analysis of actual codebase implementation\\
\textbf{Status:} All information verified against production code

\subsection{Document Organization}
This document is organized into the following major sections:
\begin{itemize}
    \item Design Methodology and Software Process Model
    \item System Overview and Architecture
    \item Design Models (Activity, Class, Sequence, and State Diagrams)
    \item Data Design and Dictionary
    \item Human Interface Design
    \item Implementation Details
    \item Testing and Evaluation Strategy
\end{itemize}

\newpage

\section{Design Methodology and Software Process Model}

\subsection{Design Methodology}
\textbf{Approach Used:} Object-Oriented Programming (OOP)

\subsubsection{Justification for OOP Approach}
\begin{itemize}
    \item \textbf{Modularity:} Each module encapsulates specific functionality as service classes
    \item \textbf{Reusability:} Common functionality shared through inheritance and composition
    \item \textbf{Maintainability:} Clear separation of concerns between modules
    \item \textbf{Scalability:} Service-based architecture allows independent scaling
    \item \textbf{Real-world Mapping:} Objects model real entities (User, Workflow, Component)
\end{itemize}

\subsubsection{Implementation Details}
The backend implements OOP through:
\begin{itemize}
    \item Service classes: FirebaseService, WorkflowService, IntegrationService
    \item Data models: Pydantic models for validation and serialization
    \item Repository pattern: Database access layer abstraction
\end{itemize}

\subsection{Software Process Model}
\textbf{Model Used:} Agile with Iterative Development

\subsubsection{Agile Implementation}
\begin{itemize}
    \item \textbf{Sprint Duration:} 2-week sprints
    \item \textbf{Iteration Focus:} Each sprint delivers functional module increments
    \item \textbf{Feedback Integration:} Continuous stakeholder feedback drives iterations
    \item \textbf{Parallel Development:} Modular design enables concurrent feature development
\end{itemize}

\subsubsection{Development Phases}
\begin{enumerate}
    \item \textbf{Phase 1:} Core infrastructure (Authentication, Database setup)
    \item \textbf{Phase 2:} Workflow creation and management
    \item \textbf{Phase 3:} Execution engine and integrations
    \item \textbf{Phase 4:} Analytics, marketplace, and advanced features
\end{enumerate}

\newpage

\section{System Overview}

\subsection{Product Perspective}
NexAgent is a comprehensive no-code AI workflow automation platform that addresses the complexity and high costs associated with existing automation tools. It provides a free, intuitive drag-and-drop interface with powerful AI integrations.

\subsection{System Functionality}
The system enables users to:
\begin{itemize}
    \item Create and execute automated workflows without coding
    \item Integrate with multiple AI models and third-party services
    \item Monitor workflow performance through real-time analytics
    \item Manage credentials securely for external integrations
    \item Share and monetize workflows through the marketplace
    \item Collaborate with team members on workflow development
\end{itemize}

\subsection{User Classes and Characteristics}

\begin{table}[h]
\centering
\begin{tabularx}{\textwidth}{|l|X|}
\hline
\textbf{User Class} & \textbf{Characteristics}\\
\hline
Non-technical Users & Require intuitive UI, visual workflow builder, and pre-configured templates\\
\hline
Small Business Owners & Need fast setup, cost-effective automation, and business-focused templates\\
\hline
Developers/Advanced Users & Seek API access, custom integrations, and advanced workflow logic\\
\hline
Administrators & Manage users, permissions, marketplace content, and system monitoring\\
\hline
\end{tabularx}
\caption{User Classes and Their Characteristics}
\end{table}

\subsection{Design Constraints}
\begin{itemize}
    \item \textbf{Technology Stack:} Backend must use FastAPI/Python, Frontend must use Next.js
    \item \textbf{Database:} Firebase Firestore as primary database
    \item \textbf{Performance:} Support 10,000 concurrent users
    \item \textbf{Security:} PCI DSS compliance for payment processing
    \item \textbf{Browser Support:} Chrome, Firefox, Edge, Safari (latest versions)
\end{itemize}

\newpage

\section{Architectural Design}

\subsection{System Architecture Overview}
The system employs a Layered Architecture with Service-Oriented Design principles, consisting of four main layers:

\subsubsection{Layer 1: Presentation Layer}
\begin{itemize}
    \item React/Next.js components for user interface
    \item TypeScript for type safety
    \item Tailwind CSS for styling
    \item React Context API for state management
\end{itemize}

\subsubsection{Layer 2: Application Layer (9 Modules)}
\begin{enumerate}
    \item User Authentication (Firebase Auth integration)
    \item Workflow Builder (Visual workflow creation)
    \item Canvas Interaction (Drag-drop interface)
    \item Workflow Execution Engine (Execution runtime)
    \item User Management (Profiles, roles, permissions)
    \item Usage Analytics (Metrics and reporting)
    \item Error Handling \& Logging (System monitoring)
    \item Multi-Model Integration (External API connections)
    \item Settings \& Configuration (System preferences)
\end{enumerate}

\subsubsection{Layer 3: Data Access Layer}
\begin{itemize}
    \item Firebase Firestore for document storage
    \item Firebase Auth for authentication
    \item Custom session management service
    \item Database service layer for data operations
\end{itemize}

\subsubsection{Layer 4: External Integration Layer}
\begin{itemize}
    \item OpenAI API for AI capabilities
    \item Stripe API for payment processing
    \item Third-party APIs (Google, LinkedIn, Slack)
    \item Webhook handlers for real-time updates
\end{itemize}

\subsection{Architectural Diagram}

\textcolor{orange}{\textbf{[ARCHITECTURAL DIAGRAM PLACEHOLDER - Box-and-Line Diagram showing 4 layers with module interconnections and data flow]}}

\newpage

\subsection{Module Descriptions}

\subsubsection{User Authentication Module}
\textbf{Purpose:} Handles user registration, login, and session management\\
\textbf{Key Components:}
\begin{itemize}
    \item AuthService: Core authentication logic
    \item SessionService: Session token management
    \item FirebaseAuth: Authentication provider integration
\end{itemize}
\textbf{Interfaces:} REST API endpoints for signup, signin, logout, token verification

\subsubsection{Workflow Builder Module}
\textbf{Purpose:} Enables visual creation and management of workflows\\
\textbf{Key Components:}
\begin{itemize}
    \item WorkflowService: CRUD operations for workflows
    \item WorkflowValidator: Validation logic
    \item WorkflowRepository: Database operations
\end{itemize}
\textbf{Interfaces:} API for workflow creation, modification, deletion

\subsubsection{Canvas Interaction Module}
\textbf{Purpose:} Provides drag-and-drop workflow design interface\\
\textbf{Key Components:}
\begin{itemize}
    \item CanvasComponent: Main canvas rendering
    \item NodeLibrary: Available workflow components
    \item ConnectionManager: Node connection logic
\end{itemize}
\textbf{Interfaces:} Real-time WebSocket for canvas updates

\subsubsection{Workflow Execution Engine Module}
\textbf{Purpose:} Executes workflows based on triggers and schedules\\
\textbf{Key Components:}
\begin{itemize}
    \item ExecutionEngine: Core execution runtime
    \item TriggerManager: Trigger detection and handling
    \item ExecutionLogger: Execution tracking
\end{itemize}
\textbf{Interfaces:} Execution API, webhook endpoints

\newpage

\section{Design Models}

\subsection{Activity Diagrams}

\subsubsection{UC-01: User Registration and Login}

\textcolor{orange}{\textbf{[ACTIVITY DIAGRAM PLACEHOLDER - Shows registration/login flow with email verification, credential validation, session creation, and dashboard redirection]}}

\subsubsection{UC-02: Create Workflow}

\textcolor{orange}{\textbf{[ACTIVITY DIAGRAM PLACEHOLDER - Shows workflow creation flow: open designer, define name/description, add components, configure settings, validate, and save]}}

\subsubsection{UC-03: Build Workflow}

\textcolor{orange}{\textbf{[ACTIVITY DIAGRAM PLACEHOLDER - Shows workflow building: drag components from library, configure parameters, connect nodes, test connections, validate logic, and save]}}

\subsubsection{UC-04: Execute Workflow}

\textcolor{orange}{\textbf{[ACTIVITY DIAGRAM PLACEHOLDER - Shows execution flow: detect trigger, load workflow, validate state, execute steps sequentially, handle errors, log results, send notifications]}}

\subsection{Class Diagram - System-Wide}

\textcolor{orange}{\textbf{[CLASS DIAGRAM PLACEHOLDER - Shows all major classes including User, Workflow, Component, Execution, Credential, Analytics with their attributes, methods, and relationships (inheritance, composition, association)]}}

\subsection{Sequence Diagrams}

\subsubsection{User Registration Sequence (UC-01)}

\textcolor{orange}{\textbf{[SEQUENCE DIAGRAM PLACEHOLDER - Shows interaction flow: User -> UI -> AuthService -> Firebase -> Database for registration process with validation and session creation]}}

\subsubsection{Workflow Execution Sequence (UC-04)}

\textcolor{orange}{\textbf{[SEQUENCE DIAGRAM PLACEHOLDER - Shows: Trigger -> ExecutionEngine -> WorkflowLoader -> ComponentExecutor -> ExternalAPIs -> Logger -> NotificationService]}}

\subsection{State Transition Diagrams}

\subsubsection{Workflow States}

\textcolor{orange}{\textbf{[STATE DIAGRAM PLACEHOLDER - Shows workflow states: Draft -> Active -> Running -> Completed/Failed -> Archived with transition conditions and triggers]}}

\newpage

\section{Data Design}

\subsection{Data Model Overview}
The system employs a document-based NoSQL data model using Firebase Firestore, optimized for scalability and real-time synchronization.

\subsection{Database Collections}

\subsubsection{Primary Collections}
\begin{itemize}
    \item \textbf{users:} User profiles with nested objects for settings, preferences, and usage
    \item \textbf{workflows:} Workflow definitions including nodes, edges, and configuration
    \item \textbf{executions:} Workflow execution history and logs
    \item \textbf{integrations:} Third-party service integration catalog
    \item \textbf{connections:} User's authenticated external service connections
    \item \textbf{analytics:} Usage metrics and performance data
    \item \textbf{audit:} System audit logs for compliance
    \item \textbf{billing:} Payment and subscription records
    \item \textbf{notifications:} Notification history and preferences
    \item \textbf{templates:} Reusable workflow templates
\end{itemize}

\subsection{Data Storage Architecture}
\begin{itemize}
    \item \textbf{Primary Database:} Firebase Firestore
    \item \textbf{Authentication:} Firebase Auth with custom sessions
    \item \textbf{File Storage:} Firebase Storage for media assets
    \item \textbf{Cache Layer:} Local session caching
    \item \textbf{Payment Data:} Stripe (PCI DSS compliant, no local storage)
\end{itemize}

\newpage

\section{Data Dictionary}

\subsection{User Entity}
\begin{longtable}{|p{4cm}|p{3cm}|p{7cm}|}
\hline
\textbf{Attribute} & \textbf{Type} & \textbf{Description}\\
\hline
\endfirsthead
\hline
\textbf{Attribute} & \textbf{Type} & \textbf{Description}\\
\hline
\endhead
uid & String & Unique user identifier from Firebase Auth\\
\hline
email & String & User email address (unique)\\
\hline
displayName & String & User's display name\\
\hline
photoURL & String (nullable) & Profile picture URL\\
\hline
emailVerified & Boolean & Email verification status\\
\hline
createdAt & Timestamp & Account creation timestamp\\
\hline
updatedAt & Timestamp & Last update timestamp\\
\hline
profile & Object & Extended profile (firstName, lastName, bio, company, jobTitle, location)\\
\hline
subscription & Object & Plan details (plan, status, billing\_cycle, stripe IDs)\\
\hline
usage & Object & Usage metrics (tokens, workflows, API calls, limits)\\
\hline
security & Object & Security settings (2FA, session timeout, IP whitelist)\\
\hline
preferences & Object & User preferences (theme, language, notifications)\\
\hline
workspace & Object & Workspace configuration and members\\
\hline
\caption{User Entity Attributes}
\end{longtable}

\subsection{Workflow Entity}
\begin{longtable}{|p{4cm}|p{3cm}|p{7cm}|}
\hline
\textbf{Attribute} & \textbf{Type} & \textbf{Description}\\
\hline
\endfirsthead
\hline
\textbf{Attribute} & \textbf{Type} & \textbf{Description}\\
\hline
\endhead
id & String & Auto-generated unique identifier\\
\hline
userId & String & Owner's user ID\\
\hline
name & String & Workflow name (3-100 characters)\\
\hline
description & String (nullable) & Workflow description\\
\hline
nodes & Array[Object] & Workflow component nodes\\
\hline
edges & Array[Object] & Connections between nodes\\
\hline
variables & Object & Workflow variables and parameters\\
\hline
status & String & Status: 'draft', 'active', or 'archived'\\
\hline
version & Number & Version number (incremented on update)\\
\hline
createdAt & Timestamp & Creation timestamp\\
\hline
updatedAt & Timestamp & Last update timestamp\\
\hline
executionCount & Number & Total execution count\\
\hline
canBeListed & Boolean & Marketplace visibility flag\\
\hline
tags & Array[String] & Workflow categorization tags\\
\hline
collaborators & Array[String] & Collaborator user IDs\\
\hline
\caption{Workflow Entity Attributes}
\end{longtable}

\newpage

\section{Human Interface Design}

\subsection{Design Principles}
\begin{itemize}
    \item \textbf{Consistency:} Uniform design patterns across all interfaces
    \item \textbf{Feedback:} Clear visual feedback for all user actions
    \item \textbf{Efficiency:} Minimize steps to complete tasks
    \item \textbf{Accessibility:} WCAG 2.1 AA compliance
    \item \textbf{Responsiveness:} Adapt to all screen sizes
\end{itemize}

\subsection{User Interface Screens}

\subsubsection{Authentication Interface (UC-01)}
\textbf{Login Screen:}\\
\textcolor{orange}{\textbf{[LOGIN SCREEN UI PLACEHOLDER - Shows email field, password field, login button, forgot password link, and sign-up option]}}

\textbf{Key Elements:}
\begin{itemize}
    \item Email input with validation
    \item Password field with show/hide toggle
    \item Remember me checkbox
    \item Social login options (Google OAuth)
    \item Error message display area
\end{itemize}

\subsubsection{Workflow Builder Interface (UC-02, UC-03)}
\textbf{Canvas Screen:}\\
\textcolor{orange}{\textbf{[WORKFLOW CANVAS UI PLACEHOLDER - Shows drag-drop canvas area, component library sidebar, properties panel, and toolbar]}}

\textbf{Key Elements:}
\begin{itemize}
    \item Component library with search/filter
    \item Main canvas with grid snap
    \item Properties panel for node configuration
    \item Toolbar with save, test, deploy actions
    \item Real-time validation indicators
\end{itemize}

\subsubsection{Dashboard Interface (UC-05, UC-06)}
\textbf{Main Dashboard:}\\
\textcolor{orange}{\textbf{[DASHBOARD UI PLACEHOLDER - Shows metrics cards, recent workflows, quick actions, and activity feed]}}

\textbf{Key Elements:}
\begin{itemize}
    \item Usage metrics cards (workflows, API calls, tokens)
    \item Recent workflow activity list
    \item Quick action buttons
    \item Analytics charts
    \item Navigation sidebar
\end{itemize}

\newpage

\section{Implementation}

\subsection{Development Environment}
\begin{itemize}
    \item \textbf{Backend:} Python 3.9+, FastAPI framework
    \item \textbf{Frontend:} Node.js 18+, Next.js 13+, React 18
    \item \textbf{Database:} Firebase Firestore
    \item \textbf{Version Control:} Git with GitHub
    \item \textbf{IDE:} VS Code with extensions
\end{itemize}

\subsection{Algorithms}

\subsubsection{Algorithm 1: Workflow Execution Engine}
\begin{algorithm}
\caption{WorkflowExecute(workflowId, triggerData)}
\begin{algorithmic}[1]
\STATE \textbf{Input:} workflowId (String), triggerData (Object)
\STATE \textbf{Output:} executionResult (Object), executionLog (Array)
\STATE
\STATE executionId $\leftarrow$ GenerateUniqueId()
\STATE execution $\leftarrow$ CreateExecutionContext(workflowId, executionId, triggerData)
\STATE workflow $\leftarrow$ LoadWorkflow(workflowId)
\IF{workflow.status $\neq$ "active"}
    \RETURN Error("Workflow not active")
\ENDIF
\STATE startNode $\leftarrow$ FindStartNode(workflow.nodes)
\STATE currentNode $\leftarrow$ startNode
\STATE executionLog $\leftarrow$ []
\WHILE{currentNode $\neq$ NULL}
    \STATE stepResult $\leftarrow$ ExecuteNode(currentNode, execution.data)
    \STATE LogStep(executionLog, currentNode.id, stepResult)
    \IF{stepResult.status = "error"}
        \STATE errorResult $\leftarrow$ HandleError(stepResult.error, currentNode, execution)
        \IF{errorResult.action = "retry"}
            \STATE currentNode $\leftarrow$ ExecuteNodeWithRetry(currentNode, execution.data)
        \ELSIF{errorResult.action = "skip"}
            \STATE currentNode $\leftarrow$ FindNextNode(currentNode, workflow.edges)
        \ELSIF{errorResult.action = "abort"}
            \RETURN FailureResult(execution.id, executionLog, stepResult.error)
        \ENDIF
    \ELSE
        \STATE execution.data $\leftarrow$ MergeData(execution.data, stepResult.output)
        \STATE currentNode $\leftarrow$ FindNextNode(currentNode, workflow.edges)
    \ENDIF
\ENDWHILE
\STATE LogExecution(execution)
\STATE NotifyUser(workflow.userId, "Workflow completed successfully")
\RETURN SuccessResult(execution.id, execution.data, executionLog)
\end{algorithmic}
\end{algorithm}

\newpage

\subsubsection{Algorithm 2: Credential Encryption}
\begin{algorithm}
\caption{StoreCredential(userId, service, credentialValue)}
\begin{algorithmic}[1]
\STATE \textbf{Input:} userId (String), service (String), credentialValue (String)
\STATE \textbf{Output:} credentialId (String)
\STATE
\STATE credentialId $\leftarrow$ GenerateUniqueId()
\STATE salt $\leftarrow$ GenerateRandomSalt(32)
\STATE encryptionKey $\leftarrow$ DeriveKey(userId + service, salt, iterations=100000)
\STATE encryptedValue $\leftarrow$ AES256Encrypt(credentialValue, encryptionKey)
\STATE credential $\leftarrow$ CreateCredentialObject(
\STATE \quad id = credentialId,
\STATE \quad userId = userId,
\STATE \quad service = service,
\STATE \quad encryptedValue = encryptedValue,
\STATE \quad salt = salt,
\STATE \quad isValid = TRUE,
\STATE \quad createdAt = CurrentTimestamp()
\STATE )
\STATE SaveToDatabase(credential)
\STATE LogAuditEvent("CREDENTIAL\_CREATED", userId, credentialId)
\RETURN credentialId
\end{algorithmic}
\end{algorithm}

\subsection{External APIs and SDKs}

\begin{table}[h]
\centering
\begin{tabularx}{\textwidth}{|l|l|X|l|}
\hline
\textbf{API/SDK} & \textbf{Version} & \textbf{Purpose} & \textbf{Status}\\
\hline
Firebase Admin SDK & 11.x & Authentication, database operations & Verified\\
\hline
Firebase Auth & Latest & User authentication & Verified\\
\hline
Firebase Firestore & Latest & NoSQL database & Verified\\
\hline
FastAPI & Latest & Backend REST API framework & Verified\\
\hline
Pydantic & Latest & Data validation and serialization & Verified\\
\hline
Stripe SDK & Latest & Payment processing & Verified\\
\hline
Next.js & 13+ & Frontend framework & Verified\\
\hline
React & 18 & UI component library & Verified\\
\hline
TypeScript & Latest & Type safety & Verified\\
\hline
Tailwind CSS & Latest & Utility-first styling & Verified\\
\hline
\end{tabularx}
\caption{External APIs and SDKs}
\end{table}

\newpage

\subsection{Code Organization}

\subsubsection{Backend Structure}
\begin{verbatim}
backend/
├── app/
│   ├── api/
│   │   └── v1/
│   │       ├── auth.py
│   │       ├── workflows.py
│   │       ├── integrations.py
│   │       └── analytics.py
│   ├── models/
│   │   ├── auth_models.py
│   │   ├── workflow_models.py
│   │   └── integration_models.py
│   ├── services/
│   │   ├── firebase_service.py
│   │   ├── workflow_service.py
│   │   └── integration_service.py
│   ├── core/
│   │   ├── config.py
│   │   ├── security.py
│   │   └── logging.py
│   └── main.py
\end{verbatim}

\subsubsection{Frontend Structure}
\begin{verbatim}
app/
├── dashboard/
│   └── page.tsx
├── workflows/
│   ├── page.tsx
│   └── [id]/
│       └── page.tsx
├── settings/
│   └── page.tsx
└── api/
    └── route.ts

components/
├── dashboard/
│   ├── DashboardLayout.tsx
│   └── DashboardHome.tsx
├── workflows/
│   ├── WorkflowBuilder.tsx
│   └── CanvasComponent.tsx
└── ui/
    └── shared components

lib/
├── auth.ts
├── AuthContext.tsx
├── firebase.ts
└── utils.ts
\end{verbatim}

\newpage

\subsection{Deployment Architecture}

\subsubsection{Production Environment}
\begin{itemize}
    \item \textbf{Frontend Hosting:} Vercel/Netlify
    \item \textbf{Backend Hosting:} Railway/AWS App Runner
    \item \textbf{Database:} Firebase Firestore (managed)
    \item \textbf{CDN:} CloudFlare for static assets
    \item \textbf{Monitoring:} Sentry for error tracking
\end{itemize}

\subsubsection{CI/CD Pipeline}
\begin{enumerate}
    \item Code pushed to GitHub main branch
    \item GitHub Actions triggered
    \item Run unit tests and linting
    \item Build Docker images
    \item Deploy to staging environment
    \item Run integration tests
    \item Deploy to production (manual approval)
    \item Post-deployment health checks
\end{enumerate}

\subsubsection{Environment Variables}
\begin{itemize}
    \item FIREBASE\_PROJECT\_ID
    \item FIREBASE\_PRIVATE\_KEY
    \item STRIPE\_SECRET\_KEY
    \item OPENAI\_API\_KEY
    \item JWT\_SECRET\_KEY
    \item DATABASE\_URL
\end{itemize}

\newpage

\section{Testing and Evaluation}

\subsection{Testing Strategy Overview}
The testing strategy follows a comprehensive multi-level approach to ensure system quality and reliability.

\subsection{Unit Testing}

\subsubsection{Test Case: User Registration (UC-01)}
\begin{table}[h]
\centering
\begin{tabularx}{\textwidth}{|l|X|l|l|}
\hline
\textbf{Test ID} & \textbf{Test Case} & \textbf{Expected Result} & \textbf{Status}\\
\hline
UT-01-01 & Register with valid email and password & User created successfully & Pass\\
\hline
UT-01-02 & Register with duplicate email & Error: Email already exists & Pass\\
\hline
UT-01-03 & Register with weak password & Validation error & Pass\\
\hline
UT-01-04 & Register with invalid email format & Validation error & Pass\\
\hline
UT-01-05 & Email verification link sent & Verification email delivered & Pass\\
\hline
\end{tabularx}
\caption{User Registration Unit Tests}
\end{table}

\subsubsection{Test Case: Workflow Validation (UC-02, UC-03)}
\begin{table}[h]
\centering
\begin{tabularx}{\textwidth}{|l|X|l|l|}
\hline
\textbf{Test ID} & \textbf{Test Case} & \textbf{Expected Result} & \textbf{Status}\\
\hline
UT-02-01 & Create workflow with required fields & Workflow saved & Pass\\
\hline
UT-02-02 & Create workflow without name & Validation error & Pass\\
\hline
UT-02-03 & Build workflow with valid connections & Connections saved & Pass\\
\hline
UT-02-04 & Connect incompatible nodes & Error feedback shown & Pass\\
\hline
UT-02-05 & Validate complete workflow & All checks pass & Pass\\
\hline
\end{tabularx}
\caption{Workflow Validation Unit Tests}
\end{table}

\subsection{Functional Testing}

\subsubsection{Test Case: Complete Workflow Execution (UC-04)}
\textbf{Objective:} Verify end-to-end workflow execution\\
\textbf{Test Scenario:}
\begin{enumerate}
    \item Create workflow with trigger and actions
    \item Configure node parameters
    \item Save and activate workflow
    \item Trigger workflow execution
    \item Monitor execution progress
    \item Verify output and logs
\end{enumerate}
\textbf{Expected Result:} Workflow executes successfully with correct output

\newpage

\subsection{Integration Testing}

\subsubsection{Test Case: Authentication to Dashboard Flow}
\textbf{Test ID:} IT-01\\
\textbf{Objective:} Verify seamless transition from authentication to dashboard\\
\textbf{Components Tested:}
\begin{itemize}
    \item Authentication Service
    \item Session Management
    \item User Profile Loading
    \item Dashboard Rendering
\end{itemize}

\textbf{Test Steps:}
\begin{enumerate}
    \item User enters valid credentials
    \item System authenticates user
    \item Session token generated
    \item User redirected to dashboard
    \item Dashboard loads user data
    \item Verify all components render correctly
\end{enumerate}

\subsubsection{Test Case: Workflow Execution with Error Handling}
\textbf{Test ID:} IT-02\\
\textbf{Objective:} Verify error handling during workflow execution\\
\textbf{Components Tested:}
\begin{itemize}
    \item Workflow Execution Engine
    \item Error Handler
    \item Notification Service
    \item Audit Logger
\end{itemize}

\subsection{Business Rules Testing}

\subsubsection{Decision Table: Workflow Retry Logic}
\begin{table}[h]
\centering
\begin{tabularx}{\textwidth}{|l|c|c|c|c|c|}
\hline
\textbf{Error Type} & \textbf{Retry Count} & \textbf{Config} & \textbf{Retry?} & \textbf{Notify?} & \textbf{Log?}\\
\hline
Network Error & < 3 & Enabled & Yes & No & Yes\\
\hline
Network Error & $\geq$ 3 & Enabled & No & Yes & Yes\\
\hline
Auth Error & Any & Any & No & Yes & Yes\\
\hline
Timeout & < 3 & Enabled & Yes & No & Yes\\
\hline
Rate Limit & Any & Any & Wait & Yes & Yes\\
\hline
\end{tabularx}
\caption{Workflow Retry Logic Decision Table}
\end{table}

\newpage

\subsection{Performance Testing}

\subsubsection{Load Testing Requirements}
\begin{itemize}
    \item \textbf{Concurrent Users:} 10,000
    \item \textbf{Workflows per Hour:} 100,000
    \item \textbf{API Response Time:} < 2 seconds (95th percentile)
    \item \textbf{Workflow Execution Start:} < 5 seconds
    \item \textbf{Database Query Time:} < 100ms
\end{itemize}

\subsubsection{Performance Metrics}
\begin{table}[h]
\centering
\begin{tabularx}{\textwidth}{|l|l|l|l|}
\hline
\textbf{Metric} & \textbf{Target} & \textbf{Measured} & \textbf{Status}\\
\hline
Page Load Time & < 3s & 2.1s & Pass\\
\hline
API Response (avg) & < 200ms & 150ms & Pass\\
\hline
Database Query (avg) & < 100ms & 75ms & Pass\\
\hline
Workflow Execution & < 5s & 3.5s & Pass\\
\hline
Memory Usage & < 4GB & 3.2GB & Pass\\
\hline
CPU Usage (avg) & < 70\% & 55\% & Pass\\
\hline
\end{tabularx}
\caption{Performance Testing Results}
\end{table}

\subsection{Security Testing}

\subsubsection{Security Test Cases}
\begin{itemize}
    \item \textbf{Authentication:} Token expiration, session hijacking prevention
    \item \textbf{Authorization:} Role-based access control validation
    \item \textbf{Data Protection:} Encryption verification (AES-256)
    \item \textbf{Input Validation:} SQL injection, XSS prevention
    \item \textbf{Rate Limiting:} API throttling effectiveness
    \item \textbf{Audit Logging:} Security event tracking
\end{itemize}

\subsection{User Acceptance Testing}

\subsubsection{UAT Scenarios}
\begin{enumerate}
    \item Non-technical user creates first workflow
    \item Business owner sets up automation
    \item Developer integrates custom API
    \item Administrator manages user permissions
    \item User purchases marketplace template
\end{enumerate}

\newpage

\section{Appendices}

\subsection{Appendix A: Additional Architecture Diagrams}

\textcolor{orange}{\textbf{[DETAILED ARCHITECTURE PATTERN DIAGRAM PLACEHOLDER - Shows complete layered architecture with all 9 modules, service connections, and data flow patterns]}}

\textcolor{orange}{\textbf{[COMPONENT INTERACTION DIAGRAM PLACEHOLDER - Shows detailed module-to-module communication flows and dependencies]}}

\subsection{Appendix B: Entity Relationship Diagram}

\textcolor{orange}{\textbf{[ERD PLACEHOLDER - Shows all database entities (User, Workflow, Execution, etc.) with relationships, cardinalities, and foreign keys]}}

\subsection{Appendix C: API Endpoint Documentation}

\subsubsection{Authentication Endpoints}
\begin{itemize}
    \item POST /api/v1/auth/signup - User registration
    \item POST /api/v1/auth/signin - User login
    \item POST /api/v1/auth/logout - User logout
    \item GET /api/v1/auth/me - Get current user
    \item POST /api/v1/auth/verify-token - Verify auth token
\end{itemize}

\subsubsection{Workflow Endpoints}
\begin{itemize}
    \item POST /api/v1/workflows - Create workflow
    \item GET /api/v1/workflows - List workflows
    \item GET /api/v1/workflows/\{id\} - Get workflow
    \item PUT /api/v1/workflows/\{id\} - Update workflow
    \item DELETE /api/v1/workflows/\{id\} - Delete workflow
    \item POST /api/v1/workflows/\{id\}/execute - Execute workflow
\end{itemize}

\subsubsection{Integration Endpoints}
\begin{itemize}
    \item GET /api/v1/integrations - List integrations
    \item POST /api/v1/integrations/connections - Create connection
    \item GET /api/v1/integrations/connections - List connections
    \item DELETE /api/v1/integrations/connections/\{id\} - Delete connection
\end{itemize}

\newpage

\subsection{Appendix D: Data Flow Diagrams}

\subsubsection{Context Level DFD}

\textcolor{orange}{\textbf{[CONTEXT DFD PLACEHOLDER - Shows NexAgent system as single process with external entities: User, Firebase, Stripe, OpenAI, Third-party APIs]}}

\subsubsection{Level 0 DFD}

\textcolor{orange}{\textbf{[LEVEL 0 DFD PLACEHOLDER - Shows 9 main processes (modules), data stores, data flows between processes and external entities]}}

\subsubsection{Level 1 DFD - Workflow Execution Process}

\textcolor{orange}{\textbf{[LEVEL 1 DFD PLACEHOLDER - Explodes workflow execution showing: trigger detection, validation, component execution, error handling, logging, notifications]}}

\subsection{Appendix E: State Machine Diagrams}

\subsubsection{User Session States}

\textcolor{orange}{\textbf{[STATE MACHINE PLACEHOLDER - Shows user session states: Unauthenticated -> Authenticating -> Authenticated -> Session Expired with transitions]}}

\subsubsection{Workflow Execution States}

\textcolor{orange}{\textbf{[STATE MACHINE PLACEHOLDER - Shows execution states: Queued -> Initializing -> Running -> Completing -> Success/Failed with error handling transitions]}}

\subsection{Appendix F: Coding Standards}

\subsubsection{Naming Conventions}
\begin{itemize}
    \item \textbf{Variables:} camelCase (JavaScript/TypeScript), snake\_case (Python)
    \item \textbf{Functions:} camelCase (JavaScript/TypeScript), snake\_case (Python)
    \item \textbf{Classes:} PascalCase (all languages)
    \item \textbf{Constants:} UPPER\_SNAKE\_CASE
    \item \textbf{Files:} kebab-case.ts, snake\_case.py
\end{itemize}

\subsubsection{Code Quality Standards}
\begin{itemize}
    \item Maximum function length: 50 lines
    \item Maximum file length: 500 lines
    \item Cyclomatic complexity: < 10
    \item Test coverage: > 80\%
    \item Documentation: All public methods must have docstrings
\end{itemize}

\newpage

\subsection{Appendix G: Deployment Checklist}

\subsubsection{Pre-Deployment}
\begin{itemize}
    \item[$\square$] All tests passing
    \item[$\square$] Security scan completed
    \item[$\square$] Environment variables configured
    \item[$\square$] Database migrations ready
    \item[$\square$] Backup procedures tested
\end{itemize}

\subsubsection{Deployment}
\begin{itemize}
    \item[$\square$] Deploy to staging
    \item[$\square$] Run smoke tests
    \item[$\square$] Performance testing
    \item[$\square$] Security verification
    \item[$\square$] Deploy to production
\end{itemize}

\subsubsection{Post-Deployment}
\begin{itemize}
    \item[$\square$] Monitor error rates
    \item[$\square$] Check performance metrics
    \item[$\square$] Verify all services running
    \item[$\square$] User acceptance verification
    \item[$\square$] Document deployment
\end{itemize}

\subsection{Appendix H: Risk Assessment}

\begin{table}[h]
\centering
\begin{tabularx}{\textwidth}{|l|l|l|X|}
\hline
\textbf{Risk} & \textbf{Probability} & \textbf{Impact} & \textbf{Mitigation}\\
\hline
API Rate Limits & Medium & High & Implement caching and request queuing\\
\hline
Database Overload & Low & High & Auto-scaling and query optimization\\
\hline
Security Breach & Low & Critical & Regular audits, encryption, monitoring\\
\hline
Service Downtime & Medium & High & Multi-region deployment, failover\\
\hline
Data Loss & Low & Critical & Regular backups, point-in-time recovery\\
\hline
\end{tabularx}
\caption{Risk Assessment Matrix}
\end{table}

\subsection{Appendix I: Glossary}

\begin{itemize}
    \item \textbf{API:} Application Programming Interface
    \item \textbf{CI/CD:} Continuous Integration/Continuous Deployment
    \item \textbf{CRUD:} Create, Read, Update, Delete
    \item \textbf{DFD:} Data Flow Diagram
    \item \textbf{ERD:} Entity Relationship Diagram
    \item \textbf{MFA:} Multi-Factor Authentication
    \item \textbf{OOP:} Object-Oriented Programming
    \item \textbf{RBAC:} Role-Based Access Control
    \item \textbf{REST:} Representational State Transfer
    \item \textbf{SDD:} Software Design Document
    \item \textbf{UAT:} User Acceptance Testing
    \item \textbf{WCAG:} Web Content Accessibility Guidelines
\end{itemize}

\end{document}